\documentclass[11pt,a4paper]{article}

\usepackage[T1]{fontenc}
\usepackage[utf8]{inputenc}
\usepackage{lmodern}
\usepackage{amsmath}
\usepackage[margin=2.4cm]{geometry}
\usepackage{booktabs}
\usepackage{longtable}
\usepackage{array}
\usepackage{xcolor}
\usepackage{listings}
\usepackage{enumitem}
\usepackage{parskip}
\usepackage[hidelinks]{hyperref}

\definecolor{rule}{HTML}{C3CBD6}
\definecolor{sunk}{HTML}{EDF0F4}
\definecolor{accent}{HTML}{1A5C7A}
\definecolor{warn}{HTML}{A84C18}
\definecolor{okgreen}{HTML}{2C6449}
\definecolor{muted}{HTML}{5C6572}

\lstdefinestyle{promptstyle}{
  basicstyle=\ttfamily\footnotesize,
  breaklines=true,
  columns=fullflexible,
  keepspaces=true,
  frame=leftline,
  framerule=2pt,
  rulecolor=\color{accent},
  framesep=8pt,
  xleftmargin=10pt,
  backgroundcolor=\color{sunk},
  showstringspaces=false,
}
\lstdefinestyle{codestyle}{
  basicstyle=\ttfamily\scriptsize,
  breaklines=true,
  columns=fullflexible,
  keepspaces=true,
  frame=single,
  rulecolor=\color{rule},
  framesep=6pt,
  xleftmargin=4pt,
  showstringspaces=false,
  captionpos=b,
  abovecaptionskip=4pt,
}
\lstset{style=codestyle}

\newenvironment{promptbox}{\medskip}{\medskip}
\newcommand{\code}[1]{\lstinline[basicstyle=\ttfamily\small]!#1!}
\newcommand{\sha}[1]{\texttt{#1}}
\newcommand{\good}[1]{\textcolor{okgreen}{\textbf{#1}}}
\newcommand{\bad}[1]{\textcolor{warn}{#1}}

\setlength{\LTpre}{6pt}\setlength{\LTpost}{6pt}
\setlength{\emergencystretch}{3em}

\title{\textbf{HIBF Fast Layout}\\[2pt]\large Claude Code session log}
\author{Session with Claude Opus~5 \textperiodcentered{} repository \texttt{HIBF\_fast\_Layout}}
\date{25--28 August 2026}

\begin{document}
\maketitle

\begin{abstract}
\noindent
A four-day debugging and optimisation session on the LSH-based HIBF layout generator. Eleven
prompts, five rounds of measured improvement to the algorithm, and three experiments that were
implemented, measured and reverted. The headline outcomes: a non-terminating layout computation
made to terminate; spurious lower-level IBFs reduced from 5901 to 128 on the 16384 dataset; a
RefSeq run that never finished in over 20 minutes brought to 158 seconds; every allocated-but-empty
technical bin removed; and an estimated 10\% reduction in index size from a greedy re-split pass.
Three latent correctness bugs were found and fixed along the way, one of which had been silently
mis-assigning input files to user bins in every layout the tool had ever produced.
\end{abstract}

\tableofcontents
\newpage

\section{Session parameters}

\begin{longtable}{@{}p{5.0cm}p{10.6cm}@{}}
\toprule
\textbf{Parameter} & \textbf{Value} \\
\midrule
\endhead
Model & \texttt{claude-opus-5} (Claude Opus~5), selected with \texttt{/model opus} at
  session start \\
Reasoning effort & \textbf{Extra High} \\
Harness & Claude Code, VS Code extension, \texttt{auto} mode (Bash-first tool policy) \\
Working directory & \texttt{/srv/public/seiler/develop/HIBF\_fast\_Layout} \\
Branch / base commit & \texttt{local}, starting from \sha{adb502a} (source unchanged since
  \sha{4fa995e}) \\
Machine & Linux 6.1.0-52-amd64, 112 cores, 1\,TB RAM \\
Compiler & GCC 15.2.0. \texttt{/import/GCC/15.2.0} became unavailable mid-session (autofs key
  present, backing server unresponsive); builds switched to
  \texttt{/group/ag\_abi/software/bin/g++-15}, the same version, via
  \texttt{make CXX=...} \\
Build flags & \texttt{-std=c++23 -O3 -march=native -DNDEBUG -fopenmp -flto} \\
Instrumented builds & \texttt{-O2}, no LTO, so symbols survive for \texttt{perf} \\
Datasets & \texttt{hibf\_paper\_lsh/\{1024\ldots524288\}} (simulated, ${\sim}65$\,GB each,
  \$BINS/64 genome groups of 64 near-duplicate pieces) and
  \texttt{RefSeq/28\_01\_2022} (25\,321 gzipped genomes, 7021 LSH clusters) \\
Layout invocation & \texttt{./generate\_layout <dir> 8 19 19 6767 0.05
  15:4,20:5,15:6,10:12 0.001 5 2 <out> <threads> [sketch\_cache]} \\
Threads & 32 for single runs, 24 when the ten simulated datasets ran back to back \\
Verification & chopper \texttt{display\_layout general}; a purpose-written structural checker;
  \texttt{function\_tests}; a numpy-based index-size proxy over the sketch cache \\
\bottomrule
\end{longtable}

\section{Prompt overview}

Elapsed times are wall-clock spans covering the whole turn, including tool execution, long
background runs and waiting. They are reconstructed from scratchpad file timestamps and, where a
commit closes the turn, from the commit time; they are therefore approximate except where noted.
Commit SHAs are the user's commits of the work produced by that prompt.

\begin{longtable}{@{}clp{5.0cm}rl@{}}
\toprule
\textbf{\#} & \textbf{Date} & \textbf{Prompt} & \textbf{Elapsed} & \textbf{Commit} \\
\midrule
\endhead
1  & 25 Aug & Why so many lower levels?            & ${\sim}$47\,min & --- \\
2  & 25 Aug & Analysis as shareable markdown       & ${\sim}$10\,min & \sha{1407a4f} \\
3  & 25 Aug & Summarise the analysis step          & ${\sim}$12\,min & --- \\
4  & 26 Aug & Fix everything; add the sketch cache & ${\sim}$110\,min & \sha{a0af981} \\
5  & 26 Aug & Update markdown + artifact           & ${\sim}$13\,min & \sha{820eba0} \\
6  & 27 Aug & RefSeq never finishes                & ${\sim}$78\,min & \sha{901822a} \\
7  & 27 Aug & Update artifact with RefSeq findings & ${\sim}$9\,min  & \sha{dcd8c26} \\
8  & 28 Aug & Fix the empty bins                   & ${\sim}$50\,min & \sha{8ba4344} \\
9  & 28 Aug & Use all allocated bins, or advise    & ${\sim}$28\,min & (folded into \sha{ef9e80b}) \\
10 & 28 Aug & Try the greedy split pass; TOC; steps & ${\sim}$40\,min & \sha{ef9e80b} \\
11 & 28 Aug & This document                        & ---             & --- \\
\bottomrule
\end{longtable}

\noindent
Only prompt~1's duration was measured directly during the session (from the scratchpad creation
time of 12:53:24 to the answer at approximately 13:40). The rest are derived from artifact
timestamps.

\section{Commits produced}

\begin{longtable}{@{}llp{8.2cm}c@{}}
\toprule
\textbf{SHA} & \textbf{Committed} & \textbf{Contents} & \textbf{Prompt} \\
\midrule
\endhead
\sha{1407a4f} & 25 Aug 15:30 & \texttt{LAYOUT\_LEVEL\_EXPLOSION.md}, \texttt{.gitignore},
  \texttt{Makefile} & 2 \\
\sha{a0af981} & 26 Aug 15:09 & The five defect fixes, three output-format fixes, sketch cache
  (6 files) & 4 \\
\sha{820eba0} & 26 Aug 15:22 & \texttt{LAYOUT\_LEVEL\_EXPLOSION.md} rewritten as
  diagnosis + resolution & 5 \\
\sha{901822a} & 27 Aug 17:18 & RefSeq performance fixes and two latent bugs (5 files) & 6 \\
\sha{a5a2f35} & 27 Aug 17:18 & \emph{User's own change:} \texttt{.fna} $\rightarrow$ \texttt{.gz}
  in the file filter & --- \\
\sha{dcd8c26} & 27 Aug 17:27 & \texttt{LAYOUT\_LEVEL\_EXPLOSION.md} + RefSeq section & 7 \\
\sha{8ba4344} & 28 Aug 12:26 & Bin packing (empty bins last) & 8 \\
\sha{ef9e80b} & 28 Aug 14:38 & Greedy re-split pass (+65 lines) and the documentation for
  prompts 9 and 10 & 9, 10 \\
\bottomrule
\end{longtable}

\newpage

\section{Prompt 1 --- Why does the layout keep adding levels?}
\textbf{25 August} \textperiodcentered{} elapsed ${\sim}$47\,min \textperiodcentered{} no commit
\begin{promptbox}
\begin{lstlisting}[style=promptstyle]
/model opus
Set model to claude-opus-5

The layout starts to add many lower levels at some point:
```
1024.layout      64
2048.layout      64
4096.layout      64
8192.layout      128
16384.layout     5901
32768.layout     1178
65536.layout     3014
131072.layout    4426
262144.layout    8105
524288.layout    15277
```

Can you figure out why? Too many are not good.

Here is the same count for another algorithm, which is more like what I would expect
```
1024.layout      64
2048.layout      64
4096.layout      64
8192.layout      128
16384.layout     128
32768.layout     192
65536.layout     256
131072.layout    384
262144.layout    512
524288.layout    768
```

in executables dir:

```
./generate_layout /srv/data/smehringer/data_simulated/hibf_paper_lsh/$BINS 8 19 19 6767 0.05 15:4,20:5,15:6,10:12 0.001 5 2 $BINS.output.layout 32
```
with $BINS being one of [1024,2048,...,524288]

to build, run `make -B` (in the executables dir)
\end{lstlisting}
\end{promptbox}
\subsection*{Output}

The dataset shape turned out to matter: \texttt{hibf\_paper\_lsh/\$BINS} is always the same
${\sim}65$\,GB of content cut into \texttt{\$BINS/64} genome groups of 64 near-identical pieces, so
LSH finds exactly one cluster per group and a whole group costs about as much as a single file
(for 1024: union 996\,118 against sum 63\,983\,552, a 1:64 ratio). The target shape is two levels.

Four defects were found, in decreasing order of impact.

\begin{enumerate}[leftmargin=*]
\item \textbf{The overflow sentinel becomes a layout.} When \code{refine_and_bin} never finds a
feasible \code{t_max}, \code{binning_core} returns \code{{0, 0, 0}} as
\code{(split_start, split_bins, merge_start)}
(\texttt{code/templates/fast\_construct\_binning\_core.tpp:140} and \texttt{:182}). With
\code{merge_start == 0}, \code{generate_hibf} turns \emph{every} non-empty technical bin into its
own lower-level IBF, single-sequence bins included. One such IBF costs about \code{global_bins}
spurious children. This is what produced the 5901. It is also a correctness bug: the sentinel
result is the partial binning built before the loop bailed out, so unplaced sequences are silently
dropped.

\item \textbf{\code{refine_and_bin} returns the last \code{t_max} it tried, not the best one.}
\texttt{generate\_hibf.tpp:89} breaks out at \code{it == p} before the result is ever scored, and
\code{b_res} is consulted only when the final result overflowed.

\item \textbf{The scoring metric drives \code{t_max} onto the cliff.}
\texttt{generate\_hibf.tpp:119} reads
\code{if (split_bin_amt == 0 || split_avg < merge_avg)}. When an IBF has no merge bins ---
the outcome we want --- \code{merge_average} returns 0
(\texttt{code/fast\_construct\_bin.cpp:24}), so the comparison is always false and \code{t_max}
is raised. The search is therefore guaranteed to climb until merge bins appear.

\item \textbf{One global bin count for the whole tree.} \code{global_bins}
(\texttt{generate\_hibf.tpp:168--176}) is computed once from the total user-bin count and reused
at every level, so the root can pack more sequences into a merge bin than the child IBF has
technical bins.
\end{enumerate}

Defects 2 and 3 together produce an \textbf{exact fixed point}: an IBF whose only merge bin
contains precisely the sequence set it was given, on which \code{generate_hibf} then recurses
forever.

\medskip\noindent\textbf{Evidence.} On the 1024 dataset the code does not terminate:

\begin{lstlisting}[style=codestyle]
[L1#15]   [it 4] t_max=62122583 used=64 split_bins=64 merge_bins=0   <- ideal, sits in b_res
[L1#15] END      t_max=62615619 nonempty=1 merge_start=63 -> children=1
[L2#0]  START n_seqs=64 sub_bins=64 lower=62122 upper=3976832 init_t_max=31554328   <- identical
[L3#0]  START n_seqs=64 sub_bins=64 lower=62122 upper=3976832 init_t_max=31554328   <- identical

[LEVEL  1] ibfs=27 -> tasks=12 total_seqs_in_tasks=768
[LEVEL 21] ibfs=12 -> tasks=12 total_seqs_in_tasks=768     <- still going when killed
\end{lstlisting}

\noindent Reproduced identically at \sha{9728f90}, so it predates the merge/seeding reorder.

\medskip\noindent\textbf{Where 5901 comes from.} The instrumented \texttt{-O2} build produced 604
lower-level IBFs on 16384 where the original \texttt{-O3 -flto} run produced 5901; hash-map
iteration order differs between builds. The breakdown of the 604:

\begin{longtable}{@{}p{11cm}r@{}}
\toprule
\textbf{Source} & \textbf{IBFs} \\
\midrule
Level 1, legitimate & 128 \\
Level 2, from the one IBF that returned the overflow sentinel (defect 1) & 128 \\
Level 2, from the 120 over-packed children (defect 4) & 348 \\
\midrule
Total & 604 \\
\bottomrule
\end{longtable}

\noindent Scaling up, $128k + {\sim}140 \approx 5901$ gives $k \approx 45$: roughly 45 of the 128
level-1 IBFs in the original run returned the overflow sentinel, against 1 here. Defect~1
dominated the reported numbers.

\medskip\noindent\textbf{Verification of the diagnosis.} A two-line scratchpad patch (score every
feasible result and keep the best; raise the search ceiling above the largest single user bin)
made 1024 terminate and cut 16384 from 604 to 367 with zero overflow sentinels --- enough to
confirm the mechanisms, not a usable fix.

\section{Prompt 2 --- Analysis as shareable markdown}
\textbf{25 August} \textperiodcentered{} elapsed ${\sim}$10\,min \textperiodcentered{} commit \sha{1407a4f}
\begin{promptbox}
\begin{lstlisting}[style=promptstyle]
Can you give me your analysis as markdown, and add the filenames to the linesnumbers, the
filenames are sometimes missing. I want to share your analysis.
\end{lstlisting}
\end{promptbox}
\subsection*{Output}

Produced \texttt{LAYOUT\_LEVEL\_EXPLOSION.md} (323 lines) in the repository root, with every source
reference written as a full \texttt{path:line} code span so it survives copy-paste, plus a
companion HTML artifact. Second-person phrasing from the chat was neutralised for sharing.

The corrected references, all verified against the source at the time:

\begin{longtable}{@{}p{9.4cm}p{5.8cm}@{}}
\toprule
\textbf{Reference} & \textbf{What it is} \\
\midrule
\endhead
\texttt{.../fast\_construct\_generate\_hibf.tpp:48} & \code{curr_upper = curr_t_max * 2}, the too-low ceiling \\
\texttt{.../fast\_construct\_generate\_hibf.tpp:89} & \code{if(it == p) break;} --- result returned unscored \\
\texttt{.../fast\_construct\_generate\_hibf.tpp:94--96} & where \code{b_res} / \code{best_t_max} are set \\
\texttt{.../fast\_construct\_generate\_hibf.tpp:119} & the \code{split_avg < merge_avg} metric \\
\texttt{.../fast\_construct\_generate\_hibf.tpp:131--133} & \code{final_res} selection \\
\texttt{.../fast\_construct\_generate\_hibf.tpp:168--176} & \code{get_sub_bins} / \code{global_bins} \\
\texttt{.../fast\_construct\_generate\_hibf.tpp:213} & the loop from \code{merge_start} that creates children \\
\texttt{.../fast\_construct\_binning\_core.tpp:140, :182} & the two \code{{0,0,0}} overflow returns \\
\texttt{.../fast\_construct\_binning\_core.tpp:142} & \code{est_sketch_size >= t_max} force-split \\
\texttt{.../fast\_construct\_binning\_core.tpp:255--259} & ``small cluster'' classification \\
\texttt{.../fast\_construct\_binning\_core.tpp:265, :271, :275} & \code{allowed_merge} truncates to 0, merge region skipped \\
\texttt{.../fast\_construct\_binning\_core.tpp:315--318, :361--363, :403} & candidates / seeding / climbing all no-op \\
\texttt{.../fast\_construct\_binning\_core.tpp:438, :58} & the fallback and \code{try_insert_sequence} \\
\texttt{code/fast\_construct\_bin.cpp:24} & \code{merge_average} returning 0 \\
\bottomrule
\end{longtable}

\noindent A caveat was flagged with the deliverable: the 604-versus-5901 discrepancy is explained
in the document, but the ``roughly 45 overflow IBFs'' figure is an inference from the reported
number, not a measurement.

\section{Prompt 3 --- Summary of the analysis step}
\textbf{25 August} \textperiodcentered{} elapsed ${\sim}$12\,min \textperiodcentered{} no commit
\begin{promptbox}
\begin{lstlisting}[style=promptstyle]
Can you give me a quick summary of the analysis step (i.e., the first prompt of this session)?
Like what model was used, how much effort, other important parameters; especially how long did it
take?
\end{lstlisting}
\end{promptbox}
\subsection*{Output}

Reconstructed from file mtimes and captured \texttt{time}/\texttt{ps} output rather than estimated.
Prompt~1 spanned 12:53:24 to approximately 13:40, about \textbf{47 minutes}.

\begin{longtable}{@{}lp{11.6cm}@{}}
\toprule
\textbf{Time} & \textbf{Activity} \\
\midrule
\endhead
12:53 & session start; read source and dataset layout \\
12:55 & \texttt{make -B} of the repository binary (31.5\,s) \\
12:55--13:05 & first 1024 run, foreground, hit the 10-minute tool timeout, killed \\
13:05--13:16 & two more 1024 runs, both still looping when killed \\
13:10 & instrumented build ready \\
13:11--13:16 & instrumented 1024 run reached \texttt{LEVEL 21}; first hard evidence of non-termination \\
13:17--13:20 & rebuilt \code{binning_core} from \sha{9728f90}, same loop, so not a regression \\
13:21--13:30 & instrumented 16384 run completed (${\sim}8.5$\,min hashing, 18\,s of layout) \\
13:33 & patched build ready \\
13:36--13:37 & patch verified on 1024 and 16384 \\
\bottomrule
\end{longtable}

\noindent The decisive accelerator was adding a sketch cache to \texttt{generate\_layout.cpp}: it
turned an ${\sim}8$-minute FASTA hashing pass into a ${\sim}2$-second load. Roughly 13 of the 47
minutes were spent before that, on two runs of the unmodified binary that could never have
finished.

\medskip\noindent Parameters reported: model \texttt{claude-opus-5}; 4 builds (1 repository build
at \texttt{-O3 -flto}, 3 instrumented at \texttt{-O2}); 9 invocations of \texttt{generate\_layout},
4 of which completed; \texttt{threads=32}; the machine sustained 1000--1500\% CPU during hashing.
Reasoning effort was \emph{not} visible to the assistant at that point and was reported as unknown
rather than guessed.
\section{Prompt 4 --- Fix everything, and make the sketch cache permanent}
\textbf{26 August} \textperiodcentered{} elapsed ${\sim}$110\,min \textperiodcentered{} commit \sha{a0af981}
\begin{promptbox}
\begin{lstlisting}[style=promptstyle]
Can you fix all the problems you found? Does it work then?
Also, I think for testing it's a good idea to also add the sketch cache permamently.
\end{lstlisting}
\end{promptbox}
\subsection*{Output}
\subsubsection*{The five defect fixes}

\begin{longtable}{@{}cp{9.0cm}p{4.6cm}@{}}
\toprule
\textbf{\#} & \textbf{Fix} & \textbf{Where} \\
\midrule
\endhead
1 & Overflow can no longer become a layout. \code{refine_and_bin} returns a scored feasible result
    or a guaranteed-valid fallback, never the \code{{0,0,0}} sentinel. &
    \texttt{generate\_hibf.tpp} \\
2 & Search ceiling raised to $\max(2 \cdot \text{init},\ \max_i(|\text{sketch}_i|/s)\cdot 17/16)$,
    so a non-overflowing \code{t_max} is always reachable. & \texttt{generate\_hibf.tpp:118, :121} \\
3 & A real objective: \code{cost = technical_bins * fullest_bin + sum of merge-bin content}. Every
    feasible candidate is scored and the best kept; the direction rule became ``overflow means
    raise, otherwise lower''. & \texttt{generate\_hibf.tpp:40, :134, :151} \\
4 & Progress guard: a merge bin may never hold its IBF's whole input, plus a
    \code{fallback_layout} that is always valid and always shrinks. Recursion depth is now finite
    by construction. & \texttt{generate\_hibf.tpp:54, :64} \\
5 & \code{try_insert_sequence} refuses to grow a bin past \code{bins} sequences, since the child
    IBF only has that many technical bins. & \texttt{binning\_core.tpp:63} \\
\bottomrule
\end{longtable}

\begin{lstlisting}[style=codestyle,caption={Fix 2: above max\_single nothing is split, so binning\_core cannot run out of bins}]
const size_t feasible_t_max = max_single + max_single/16 + 1;  // max_single = max_i(|sketch_i|/s)
size_t curr_upper = std::max(curr_t_max * 2, feasible_t_max);
\end{lstlisting}

\subsubsection*{Three output-format bugs found while validating}

None of these relate to the level explosion; they surfaced when the layouts were first checked
against chopper's parser. All three are in \texttt{code/fast\_construct\_bin.cpp}.

\begin{enumerate}[leftmargin=*]
\item \textbf{Line 124 --- the lower-level header was missing its colon.} The top-level line had
it, the lower-level line did not, so \code{layout::read_from} could not parse the header:
\begin{lstlisting}[style=codestyle]
wrote:    #LOWER_LEVEL_IBF_48 fullest_technical_bin_idx17
expected: #LOWER_LEVEL_IBF_48 fullest_technical_bin_idx:17
\end{lstlisting}
\item \textbf{Line 43 --- \code{@CHOPPER_USER_BINS_END} had no trailing newline}, so it ran into
the following \code{@CHOPPER_CONFIG} and produced the line
\code{@CHOPPER_USER_BINS_END@CHOPPER_CONFIG}.
\item \textbf{Lines 38--44 --- the user-bin section was emitted in
\code{std::unordered_map} order.} \code{chopper::layout::read_filenames_from} reads that section
\emph{positionally} (\code{filenames.emplace_back()}); the \code{@<id>} on each line is only
cross-checked under \code{assert}. In a release build every user bin was therefore silently paired
with a file belonging to a different one. \textbf{Every layout the tool produced before
\sha{a0af981} had scrambled user-bin-to-file assignments} --- relevant before comparing against
older benchmark numbers.
\end{enumerate}

\subsubsection*{The sketch cache}

\code{read_sketch_cache} / \code{write_sketch_cache} in
\texttt{code/fast\_construct\_hashing.cpp:97} and \texttt{:127}, declared in
\texttt{fast\_construct\_hashing.h}, wired to an optional 13th argument. The file header stores
\code{dir_path}, \code{q}, \code{k}, \code{w}, \code{seed} and \code{s}, and the loader refuses a
cache written with different parameters, so a stale one cannot be picked up silently.

\subsubsection*{Results}

\begin{lstlisting}[style=codestyle,caption={The shape every dataset now has}]
[hibf] level 0: 1 IBF(s), 128 merge bin(s) to expand
[hibf] level 1: 128 IBF(s), 0 merge bin(s) to expand
\end{lstlisting}

\begin{longtable}{@{}rrrrc@{}}
\toprule
\textbf{\$BINS} & \textbf{before} & \textbf{after} & \textbf{chopper} & \textbf{levels} \\
\midrule
\endhead
1024   & \bad{never terminates} & \good{16}  & 64  & 2 \\
2048   & 64                     & \good{32}  & 64  & 2 \\
4096   & 64                     & \good{64}  & 64  & 2 \\
8192   & 128                    & \good{128} & 128 & 2 \\
16384  & \bad{5901}             & \good{128} & 128 & 2 \\
32768  & \bad{1178}             & \good{192} & 192 & 2 \\
65536  & \bad{3014}             & \good{256} & 256 & 2 \\
131072 & \bad{4426}             & \good{384} & 384 & 2 \\
262144 & \bad{8105}             & \good{514} & 512 & 3 \\
524288 & \bad{15277}            & \good{785} & 768 & 3 \\
\bottomrule
\end{longtable}

\noindent Verified three ways: chopper's \texttt{display\_layout general} parses them (for 16384 it
reports 128 top-level bins, all \texttt{merged}, \texttt{ub\_count=128}, sizes 7.81e6 within 0.5\%
of each other); a structural checker confirms every user bin appears exactly once, headers and
content agree and are in level order, and technical bin spans do not overlap; and
\texttt{function\_tests} output is byte-identical to a build from before the fix (51 TRUE / 4
FALSE --- those four already failed).

\medskip\noindent\textbf{Sketch cache effect:} 16384 went from 363\,s to 13\,s, 262144 from 983\,s
to 54\,s, with a byte-identical layout (\texttt{cmp} clean).

\medskip\noindent\textbf{Residual, reported honestly:} 262144 and 524288 keep a small third level
(2 and 17 IBFs). Measured cause: the occupancy cap is exactly \code{bins}, leaving a child no slack
for a sequence that must split across two technical bins. For 262144 the geometry is exactly tight
--- every level-1 IBF holds exactly 512 user bins in exactly 512 technical bins ($512 \times 512 =
262144$) --- and lowering the cap does not help because the root has no slack either.

\section{Prompt 5 --- Update the markdown and the artifact}
\textbf{26 August} \textperiodcentered{} elapsed ${\sim}$13\,min \textperiodcentered{} commit \sha{820eba0}
\begin{promptbox}
\begin{lstlisting}[style=promptstyle]
Yes, I want LAYOUT_LEVEL_EXPLOSION.md and the artifact updated with the fixes and the results
table.
\end{lstlisting}
\end{promptbox}
\subsection*{Output}

\texttt{LAYOUT\_LEVEL\_EXPLOSION.md} was restructured from a pure diagnosis into a
diagnosis-and-resolution document (418 lines), and the artifact was redeployed to the same URL.

\begin{itemize}[leftmargin=*]
\item Result section moved to the front with the full before/after/chopper table and the
three-way validation.
\item The diagnosis was kept intact; \texttt{Cause N} was renamed \texttt{Defect N} and put in
past tense.
\item A new ``The fix'' section replaced ``Suggested fixes'', with line numbers against
\sha{a0af981}.
\item New sections on the three output-format bugs and on what remained open.
\item A ``Testing'' section documenting the sketch cache and the cold-versus-cached timings.
\end{itemize}

\noindent Two observations were flagged rather than acted on: the user's own tweak to
\code{write_linkage} in \sha{a0af981} (an indexed loop instead of the sorted one) produces identical
output, verified by re-running 16384; and the uncommitted \texttt{Makefile} change to
\texttt{CXX = g++} needs to resolve to a C++23-capable compiler, since the system \texttt{g++} on
this machine is 12.2.

\section{Prompt 6 --- RefSeq never finishes}
\textbf{27 August} \textperiodcentered{} elapsed ${\sim}$78\,min \textperiodcentered{} commit \sha{901822a}
\begin{promptbox}
\begin{lstlisting}[style=promptstyle]
Can you have a look at another dataset:
```
./generate_layout "/srv/data/datasets/RefSeq/28_01_2022/" 8 19 19 6767 0.05 15:4,20:5,15:6,10:12 0.001 5 2 refseq.output.layout 32 refseq.bin
```
The chached sketches are at `/srv/public/seiler/develop/HIBF_fast_Layout/executables/refseq.bin`
if you want to reuse them.

the problem is that it takes very long; i've terminated it after 20 minutes.
Do you see any problems that might cause this?
If so, fix them.

Another note, in the staged changes, I changed `.fna` to `.gz`. This is because this data set is
gzipped and I just need the loop to add them to the file list. seqan3's sequence file input
handles the compression (otherwise, it would throw). This workaround is fine for now.
\end{lstlisting}
\end{promptbox}
\subsection*{Output}

RefSeq is a fundamentally different shape from the simulated data, and that difference is what
exposed the problems.

\begin{longtable}{@{}p{6.6cm}rr@{}}
\toprule
 & \textbf{simulated 16384} & \textbf{RefSeq} \\
\midrule
\endhead
user bins & 16384 & 25321 \\
LSH top-level clusters & 256 & \bad{7021} \\
technical bins (\code{global_bins}) & 128 & 192 \\
clusters per technical bin & 2 & \bad{37} \\
sequence sizes & uniform & highly variable \\
\bottomrule
\end{longtable}

\noindent The simulated data is 64 near-duplicates per genome, so almost every insertion is
accepted on the first attempt. RefSeq is genuinely diverse, so almost every attempt is
\emph{rejected} --- and the rejection path is where all the cost and both latent bugs were.

\subsubsection*{Finding 1 --- a rejected insertion cost a full sketch scan, thrown away}

\code{try_insert_sequence} probed the bin's hash set once per element of the sequence's sketch,
computed the union size, and only then decided. On rejection all of that work was discarded, and
seeding, late merging and the fallback each offer the same sequence to up to \code{bins} bins in
turn, giving $O(n_{\text{seqs}} \cdot \text{bins} \cdot |\text{sketch}|)$ hash probes.

\begin{lstlisting}[style=codestyle,caption={One root binning pass, before the fix}]
[bc] nseq=25321 bins=192 t_max=151249359 clust=7021 | ... fallback=68501ms TOTAL=76044ms
     try=1391818 rej_tmax=253345 lookups=1073.8M
\end{lstlisting}

\noindent 1.07 billion hash probes in a single pass, 68\,s of it in the fallback alone, repeated
for every \code{t_max}. Fixed at \texttt{binning\_core.tpp:68--85} with two exact checks: the union
can never be smaller than either side, which rejects a full bin without touching the hash table at
all, and the loop stops as soon as the running union is known to exceed \code{t_max}. Neither
changes which insertions are accepted.

\subsubsection*{Finding 2 --- work independent of \texttt{t\_max} redone for every \texttt{t\_max}}

The estimated size of every top-level cluster (7021 \code{get_union_size_ptr} calls, 2.1\,s per
pass) and, for child IBFs, \code{filter_LSH} over the entire 33\,569-cluster forest, were both
recomputed on each of the six refinement passes. The cluster sizes are now cached per IBF
(\texttt{binning\_core.tpp:216}, cache owned at \texttt{generate\_hibf.tpp:137}) and
\code{filter_LSH} is hoisted to run once per IBF (\texttt{generate\_hibf.tpp:132}). Visible in the
trace: \texttt{top=2119ms} on the first pass, \texttt{top=1ms} on the five after it.

\subsubsection*{Finding 3 --- an unsigned underflow that spins about $2^{64}$ times}

\begin{lstlisting}[style=codestyle]
size_t num_seeding_bins = seeding_max_bin - start;   // wraps when start > seeding_max_bin
\end{lstlisting}

\noindent Late merging reserves the tail of the bin range and splitting consumes the head, so the
seeding window can be empty; the subtraction then wraps to about 1.8e19 and the loop below it
spins essentially forever on cheap operations. It fired \textbf{268 times} in one RefSeq run, e.g.
\texttt{start=192 seeding\_max\_bin=191 bins=192} --- off by a single bin. The symptom was one
child IBF at 100\% CPU while its 191 siblings each finished in 0.5\,s. Unchanged since at least
\sha{4fa995e}. Clamped to 0 at \texttt{binning\_core.tpp:396}.

\subsubsection*{Finding 4 --- an out-of-bounds read that silently dropped user bins}

\begin{lstlisting}[style=codestyle]
size_t best_b = start;
size_t best_fill = track_fill[start];   // start == bins reads past the end
...
try_insert_sequence(seq, best_b, true); // and this rejects b >= bins, so nothing is placed
\end{lstlisting}

\noindent The sequence never reached the child IBF, but it still had a root-level entry pointing at
the merge bin it came from --- so 342 user bins ended up claiming technical bins belonging to other
user bins: corrupt entries rather than missing ones. Also pre-existing. Fixed at
\texttt{binning\_core.tpp:536--537} by retrying over the bins outside the split region and
reporting overflow if splitting claimed all of them, which the \code{t_max} search then resolves.
The merge window is clamped the same way at \texttt{binning\_core.tpp:308}, because bins below
\code{merge_bins} hold pieces of split sequences and \code{record_bins} needs those to stay a
contiguous run.

\medskip\noindent A first attempt that simply searched from bin~0 traded the 342 dropped user bins
for 153 impossible four-level paths --- the force path had started writing into split bins ---
which is how that constraint surfaced.

\subsubsection*{Result}

\begin{longtable}{@{}p{7.2cm}p{4.0cm}p{3.4cm}@{}}
\toprule
 & \textbf{before} & \textbf{after} \\
\midrule
RefSeq, 25\,321 genomes & \bad{terminated after 20\,min, no output} & \good{158\,s} \\
\bottomrule
\end{longtable}

\noindent Structure: $1 + 192 + 696 = 889$ IBFs over 3 levels, all 25\,321 user bins present
exactly once, no overlapping technical bin spans. The ten simulated datasets produce identical
layouts and all got faster: 262144 went from 54\,s to 33\,s from cache, with fewer threads
(24 against 30).

\medskip\noindent\textbf{Diagnostic method note:} \texttt{perf} attribution was useless here
because LTO inlined everything into \code{binning_core}, and \texttt{gdb} could not attach
(\texttt{ptrace} blocked). The problem was localised instead by building an \texttt{-O2}
instrumented copy with per-phase timers inside \code{binning_core}, which showed all phases
completing in ${\sim}14$\,s and then nothing --- pointing directly at the seeding loop.

\section{Prompt 7 --- Update the artifact with the RefSeq findings}
\textbf{27 August} \textperiodcentered{} elapsed ${\sim}$9\,min \textperiodcentered{} commit \sha{dcd8c26}
\begin{promptbox}
\begin{lstlisting}[style=promptstyle]
Can you also update the artifact with the problem description and what you found out and how you
fixed it?

[sent mid-turn, while the RefSeq investigation of prompt 6 was still running]
\end{lstlisting}
\end{promptbox}
\subsection*{Output}

Both documents gained a ``RefSeq: a different failure'' part: the dataset comparison table, the four
findings as a ranked list, the instrumented traces, and the result. The markdown was retitled
\emph{Two failures in the fast layout, diagnosed and fixed} since it now covered two investigations,
and grew to 545 lines. The artifact was redeployed to the same URL with a matching section.
\section{Prompt 8 --- Fix the empty bins}
\textbf{28 August} \textperiodcentered{} elapsed ${\sim}$50\,min \textperiodcentered{} commit \sha{8ba4344}
\begin{promptbox}
\begin{lstlisting}[style=promptstyle]
Can you try to fix the issue with the many empty bins?
\end{lstlisting}
\end{promptbox}
\subsection*{Output}

\code{binning_core} sorted its bins ascending by size, which put the \textbf{empty bins at the low
indices} and the used ones at the top. But seqan::hibf derives an IBF's technical bin count from
the \emph{highest} index in the layout and allocates that many rounded up to a multiple of 64:

\begin{lstlisting}[style=codestyle]
// seqan::hibf/src/layout/graph.cpp:81
number_of_technical_bins = std::max(number_of_technical_bins, bin + num_tbs);

// seqan::hibf/src/interleaved_bloom_filter.cpp:41-43
bin_words      = divide_and_ceil(bins, 64u);
technical_bins = bin_words * 64u;
resize(technical_bins * bin_size_);
\end{lstlisting}

\noindent Trailing empty bins are therefore free; leading ones are paid for in full. An IBF using
16 bins at indices 48--63 still declared, and allocated, 64.

\medskip\noindent\textbf{The fix} inverts the ordering to \textbf{split bins, then merge bins, then
the empties}, packing used bins at $0 \ldots k-1$. That required replacing the two
\code{partition_point} calls with an explicit scan, since \texttt{[size 1][size >= 2][empty]} is not
monotone for either predicate. It is a pure relabelling: same bins, same contents, same sequences.

\begin{longtable}{@{}lrrr@{}}
\toprule
\textbf{RefSeq} & \textbf{declared tbins} & \textbf{empty} & \textbf{allocated (64-padded)} \\
\midrule
before & 170\,688 & \bad{41\,949 (25\%)} & 170\,688 \\
after  & 128\,739 & \good{0}             & \good{145\,792 ($-15\%$)} \\
\bottomrule
\end{longtable}

\noindent Every dataset now reports \texttt{empty\_tbins=0}, all layouts still validate, every
lower-level IBF count is unchanged, and RefSeq runtime is unchanged.

\subsubsection*{Tried and reverted: charging allocated 64-blocks in \texttt{layout\_cost}}

Arguably the more correct model, since the IBF really only allocates
$\lceil \text{used}/64 \rceil \cdot 64$ bins. Measured, it cuts RefSeq's allocated bins further to
126\,144 (waste 12\% to 7\%) but buys that with more IBFs: RefSeq 889 to 1059, and 524288 from 785
to 890. Since the net memory effect could not be confirmed without measuring real bin sizes, and
``too many lower levels'' was the original complaint, it was reverted, with the measurement recorded
in a comment at \texttt{generate\_hibf.tpp:38--43}.

\section{Prompt 9 --- Use all allocated bins, or advise}
\textbf{28 August} \textperiodcentered{} elapsed ${\sim}$28\,min \textperiodcentered{} documentation folded into \sha{ef9e80b}
\begin{promptbox}
\begin{lstlisting}[style=promptstyle]
Can you try to use all allocated bins (i.e., in your example, use 128 instead of 100, because we
allocate 128 anyway). If it's not trivial to do, or it changes the algorithm drastically, do not
do it, and just summarise possible avenues to achieve the goal.
Afterwards, please update the markdown and artifact with the changes/findings from the previous
and this prompt.
\end{lstlisting}
\end{promptbox}
\subsection*{Output}

The recommendation at this point was \textbf{not to implement it}, on the following measurements.

\begin{itemize}[leftmargin=*]
\item \textbf{The free lever does not work.} \texttt{p=12} instead of \texttt{5} leaves the padding
waste at 12\% and doubles the runtime (303\,s against 143\,s). The \code{t_max} bisection already
sits at the overflow boundary.
\item \textbf{The waste is thin and spread out}, not concentrated:
\end{itemize}

\begin{longtable}{@{}lrr@{}}
\toprule
\textbf{used bins} & \textbf{\#IBFs} & \textbf{padding wasted} \\
\midrule
1--64    &  36 &   157 \\
65--128  & 317 & 7\,507 \\
129--192 & 536 & 9\,389 \\
\bottomrule
\end{longtable}

\begin{itemize}[leftmargin=*]
\item \textbf{Best case if it worked perfectly:} a split bin's fill is exactly its sequence's sketch
size, so this is computable from the cache. Greedily handing every spare bin to the currently
fullest user bin, with fcorr applied, gives a \textbf{16.1\%} smaller \code{bin_size} across the 716
IBFs that could use it --- mostly modest (302 under 10\%, 295 between 10 and 20\%).
\end{itemize}

\noindent Getting it would need an \textbf{fcorr-aware fill model} (\code{track_fill} holds raw
content and \code{layout_cost} ignores fcorr, so the search cannot see the trade-off) and a
\textbf{re-splitting pass} preserving the split-bin contiguity invariant --- the same invariant
behind two of the earlier bugs. That was judged past the ``not trivial'' bar.

\subsubsection*{Avenues offered, best effort-to-benefit first}

\begin{enumerate}[leftmargin=*]
\item \textbf{Size \code{global_bins} per IBF instead of globally.} \code{get_sub_bins} hands every
IBF $\lceil \sqrt{N_{\text{total}}}/64 \rceil \cdot 64$; a child holding 133 sequences gets 192
bins. Worth about 12\% of allocated bins, no re-splitting. Changes tree shape, needs full
re-verification.
\item \textbf{Charge allocated 64-blocks in \code{layout_cost}.} Already written and measured (see
prompt~8); the cheapest lever if the extra IBFs turn out not to matter.
\item \textbf{Greedy re-split post-pass.} The 16.1\% option.
\item \textbf{Accept it.} 12\% of allocated technical bins on RefSeq, 0--4\% on the simulated data.
\end{enumerate}

\noindent Both documents were updated with an ``Empty technical bins'' section carrying the packing
fix, the before/after numbers and these avenues.

\section{Prompt 10 --- Try the greedy split pass; add a TOC and the steps}
\textbf{28 August} \textperiodcentered{} elapsed ${\sim}$40\,min \textperiodcentered{} commit \sha{ef9e80b}
\begin{promptbox}
\begin{lstlisting}[style=promptstyle]
Can you try out the greedy split pass? If this turns out to work well, please update the markdown
and artifact. For the artifact, please also add a toc if possible; and somehow indicate the
fixes/steps taken -- what I mean is that we iteratively improved the algorithm and I'd like an
indication of the steps, I.e. what did we do first, what was the next step, etc.
\end{lstlisting}
\end{promptbox}
\subsection*{Output}

The pass works and was kept. 65 lines in
\texttt{code/templates/fast\_construct\_binning\_core.tpp}, after placement and before the final
sort:

\begin{lstlisting}[style=codestyle]
used   = non-empty bins
target = min(bins, ceil(used / 64) * 64)      // what the HIBF allocates anyway
spare  = target - used

while spare > 0:
    take the split user bin with the largest fcorrs[k] * (content / k)
    stop if it is not larger than the fullest merge bin      (merge bins cannot be split)
    stop if giving it one more bin does not lower its demand (fcorrs[k+1] outgrows the division)
    give it a free bin, re-partition its content evenly over k+1 bins
\end{lstlisting}

\noindent Three properties keep it cheap and safe. \textbf{Merged bins are excluded} --- splitting
one would mean a whole new child IBF, so the fullest merged bin is a floor on what the pass can
achieve and it stops there. \textbf{No contiguity work is needed}: the final sort orders by
\texttt{(size, front sequence)}, so the $k+1$ technical bins of a user bin end up adjacent by
construction, which is exactly what \code{record_bins} requires. And \textbf{\code{bin_sketches} is
never touched} --- it is local to \code{binning_core} and dead after this point, so only \code{res}
and \code{track_fill} need updating.

\subsubsection*{Result}

Measured as
\[
  \sum_{\text{IBF}} \bigl(\text{allocated technical bins}\bigr) \times \bigl(\text{bin size}\bigr),
\]
where a split user bin's demand is $\text{fcorrs}[k] \cdot \text{content}/k$ and a merged bin's is
the union of its whole subtree, computed from the sketch cache with numpy.

\begin{longtable}{@{}lrrr@{}}
\toprule
\textbf{RefSeq} & \textbf{IBFs} & \textbf{allocated tbins} & \textbf{estimated index size} \\
\midrule
packing only          & 889          & 145\,792 & 100.0\% \\
+ greedy re-split     & \good{883}   & 153\,216 & \good{89.8\%} \\
+ fcorr-aware cost    & 1105         & 186\,368 & 94.3\% \\
\bottomrule
\end{longtable}

\noindent A 10\% smaller index, slightly \emph{fewer} IBFs, runtime unchanged at ${\sim}141$\,s,
padding waste 12\% to 1\%. The ten simulated datasets are unchanged --- same lower-level IBF counts,
same estimated size (1024, 16384 and 524288 all measured at exactly 100.0\%). That is the expected
profile: the pass only fires where an IBF has both spare bins and split bins.

\subsubsection*{The companion change that backfired}

Making \code{layout_cost} fcorr-aware --- replacing raw \code{max_fill} with
$\max_{\text{user bins}} \text{fcorrs}[k] \cdot (\text{content}/k)$ --- looks necessary, since the
cost model otherwise prices splitting as free. Measured, it makes things \textbf{worse}: 94.3\%
against 89.8\%, with IBFs jumping from 883 to 1105.
\code{sub_bins * max_demand} with an fcorr-inflated demand overstates the cost of exactly the
split-heavy layouts the pass is designed to produce, so the \code{t_max} search retreats to more,
smaller merged bins. Reverted.

\subsubsection*{A measurement trap worth recording}

The first size proxy summed $\text{allocated\_tbins} \times \text{bin\_size}$ over only those IBFs
with no merged bins, because a merged bin's demand needs the union of its whole subtree. That set of
IBFs differs between layouts, so the sum compared different things and reported the pass as 10\%
\emph{worse}. Computing the merged unions properly reversed the sign. This was corrected in the
session and is recorded in both documents.

\subsubsection*{Document changes requested}

The artifact gained a two-column, anchor-linked \textbf{table of contents} over all 16 sections, and
a \textbf{``How we got here''} stepped timeline of the five rounds in order --- each with its commit
and headline metric --- plus a greyed-out final entry for the three experiments that were written,
measured and reverted, since what they ruled out shaped the next step.

\section{Prompt 11 --- This document}
\textbf{28 August} \textperiodcentered{} no commit
\begin{promptbox}
\begin{lstlisting}[style=promptstyle]
Can you please create a latex document of this complete session? Store it locally in the working
directory as `claude.tex`. It should include parameters (your effort is set to `Extra High`),
prompts and output. You do not need to include the individual think steps -- the prompt and final
output for the prompt is enough. If you know how long a promp took to process, include the total
time to process a prompt. If the result of a prompt was added as commit to the repository, please
include the commit sha.
\end{lstlisting}
\end{promptbox}
\subsection*{Output}

This file, \texttt{claude.tex}, written to the repository root. Timings are reconstructed from
scratchpad artifact timestamps and commit times as described in section~2; only prompt~1's duration
was measured directly during the session. The document was compiled with \texttt{pdflatex} to
confirm it builds.

\section{Cumulative outcome}

\begin{longtable}{@{}p{6.4cm}p{4.2cm}p{4.2cm}@{}}
\toprule
\textbf{Metric} & \textbf{Session start (\sha{adb502a})} & \textbf{Session end (\sha{ef9e80b})} \\
\midrule
1024, lower-level IBFs   & \bad{never terminates} & \good{16} \\
16384, lower-level IBFs  & \bad{5901}             & \good{128} \\
32768                    & \bad{1178}             & \good{192} \\
65536                    & \bad{3014}             & \good{256} \\
131072                   & \bad{4426}             & \good{384} \\
262144                   & \bad{8105}             & \good{514} \\
524288                   & \bad{15277}            & \good{785} \\
RefSeq (25\,321 genomes) & \bad{$>$20\,min, no output} & \good{141\,s, 883 IBFs} \\
RefSeq empty technical bins & \bad{41\,949}       & \good{0} \\
RefSeq estimated index size & 100\%               & \good{89.8\%} \\
User bin to file mapping & \bad{scrambled (silent)} & \good{correct} \\
Layout parseable by chopper & \bad{no (three format bugs)} & \good{yes} \\
16384 re-run from cache  & 363\,s                 & \good{12\,s} \\
\bottomrule
\end{longtable}

\subsection*{Bugs fixed, by kind}

\begin{longtable}{@{}p{4.2cm}p{6.0cm}p{4.6cm}@{}}
\toprule
\textbf{Kind} & \textbf{Bug} & \textbf{Consequence} \\
\midrule
\endhead
Non-termination & Refinement returns the last, unscored \code{t_max}; metric climbs onto the
  collapse point & 1024 looped forever \\
Layout explosion & \code{{0,0,0}} overflow sentinel emitted as a layout & thousands of spurious
  IBFs, plus dropped sequences \\
Layout explosion & \code{global_bins} fixed tree-wide & forced merging one level down \\
Correctness & User-bin section written in hash order, read positionally & every user bin paired
  with the wrong file \\
Correctness & Missing colon in lower-level header & layout unparseable \\
Correctness & Missing newline after \code{@CHOPPER_USER_BINS_END} & layout unparseable \\
Correctness & Out-of-bounds \code{track_fill[start]} with \code{start == bins} & 342 user bins
  claimed other bins' technical bins \\
Hang & \code{seeding_max_bin - start} unsigned underflow & ${\sim}2^{64}$-iteration spin \\
Performance & Rejected insertion cost a full sketch scan & 1.07e9 hash probes per pass \\
Performance & \code{t_max}-independent work redone per \code{t_max} & 2.1\,s per pass, x6 \\
Space & Empty bins sorted first, so allocated & 41\,949 bins paid for and unused \\
Space & Allocated-but-unused bins never spent & 10\% larger index \\
\bottomrule
\end{longtable}

\subsection*{Still open}

\begin{itemize}[leftmargin=*]
\item 262144 and 524288 keep a small third level (2 and 17 IBFs); the occupancy cap is exactly
\code{bins}, and for 262144 the geometry is exactly tight so lowering it does not help.
\item RefSeq keeps a third level of 690 IBFs, and the root binning still dominates its runtime
(${\sim}150$\,s of 158\,s) --- real work now rather than waste.
\item \texttt{function\_tests} reports 4 failures (seeding-after-splitting, climbing, splitting,
merge/split isolation). These predate the session: output is byte-identical to a build from before
any change.
\item Content line order is not stable across builds, since \code{write_content} iterates an
\code{unordered_map}. The parser does not care, but sorting by user bin id would make output
byte-reproducible.
\item Whether the block-aware \code{layout_cost} is a net win is unresolved; a real
\texttt{raptor build} size comparison would settle it.
\item RefSeq's layout is validated by the structural checker but not yet by chopper's
\texttt{display\_layout}, which was stopped as too expensive.
\end{itemize}

\subsection*{Artefacts produced}

\begin{itemize}[leftmargin=*]
\item \texttt{LAYOUT\_LEVEL\_EXPLOSION.md} --- 699 lines, the full diagnosis and resolution record.
\item A published HTML artifact carrying the same material with a table of contents and a stepped
timeline.
\item \texttt{check\_layout.py} --- structural validator for chopper/raptor layout files.
\item \texttt{blocks.py}, \texttt{size\_proxy.py}, \texttt{full\_size.py} --- technical bin and
index-size analysis over the sketch cache.
\item The sketch cache itself, now a permanent feature of \texttt{generate\_layout}.
\end{itemize}

\end{document}